% =============================================================================
% CHAPITRE 1 - Introduction
% Source : consigne/revue_litterature_v2_DSR.tex (lignes 73 a 236)
% =============================================================================

\chapter{Introduction}\label{chap:introduction}

En 2025, l\textquotesingle intelligence artificielle générative a
franchi un seuil décisif dans le domaine du génie logiciel : les grands
modèles de langage (LLM) ne se contentent plus de compléter du code
ligne par ligne, ils orchestrent désormais des workflows complets de
développement, de migration et de maintenance logicielle. Chez Google,
39 migrations de code distinctes totalisant 93 574 éditions
individuelles ont été réalisées avec l\textquotesingle assistance de LLM
en seulement douze mois, avec 74,45 \% des changements appliqués sans
modification humaine (Ziftci et al., 2025). Amazon Q Developer
transforme automatiquement des applications Java 8 vers Java 17 avec 79
\% du code appliqué directement (Amazon, 2024). IBM Watsonx
s\textquotesingle attaque à la modernisation de systèmes COBOL, estimant
que 775 milliards de lignes de code legacy restent en production à
l\textquotesingle échelle mondiale (IBM, 2023).

Ces avancées industrielles convergent vers un constat : la
transformation de code à grande échelle ne relève plus de la correction
unitaire de bugs, mais de la coordination de multiples agents logiciels
opérant simultanément sur des bases de code massives, interdépendantes
et critiques pour l\textquotesingle activité des entreprises.
Les organisations font face à des choix architecturaux complexes :
certaines cherchent à casser des monolithes pour migrer vers des
microservices, d\textquotesingle autres modernisent progressivement
des composants tout en maintenant le système principal en production,
d\textquotesingle autres encore doivent arbitrer entre refonte complète
et migration incrémentale. Cette diversité de stratégies exige des
cadres d\textquotesingle orchestration flexibles, capables de
s\textquotesingle adapter à des objectifs variés plutôt que de
prescrire un workflow unique.

Or, l\textquotesingle architecture dominante pour orchestrer ces agents
repose sur des modèles hiérarchiques où un superviseur central planifie,
distribue et contrôle les tâches de chaque agent subordonné, à
l\textquotesingle image des frameworks MetaGPT (Hong et al., 2024),
CrewAI ou AutoGen. Cette approche centralisée reproduit, dans le monde
logiciel, les structures organisationnelles classiques du
command-and-control.

Pourtant, la littérature récente révèle des limitations structurelles
préoccupantes de ces architectures hiérarchiques. Le framework MAST
identifie 14 modes d\textquotesingle échec distincts des systèmes
multi-agents LLM, avec des taux d\textquotesingle échec allant de 41 \%
à 86,7 \% sur sept frameworks état de l\textquotesingle art (Cemri et
al., 2025). L\textquotesingle approche Agentless démontre
qu\textquotesingle un pipeline séquentiel simple atteint 32 \% de
résolution sur SWE-bench Lite pour seulement 0,70 USD par problème,
surpassant tous les systèmes agents open source (Xia et al., 2024). Gao
et al. (2025) observent que les gains de performance des systèmes
multi-agents chutent de 10-16 \% à seulement 0,8-3 \% lorsque des
modèles frontières sont utilisés, questionnant la valeur ajoutée réelle
de la complexité agentique.

Face à ces constats, une alternative théorique existe : la stigmergie,
mécanisme de coordination indirecte où les traces laissées par une
action dans un médium stimulent les actions subséquentes (Heylighen,
2016a). Conceptualisée par Grassé (1959) pour expliquer la construction
collective des termitières sans superviseur central, la stigmergie a été
formalisée par Bonabeau, Dorigo et Theraulaz (1999) comme fondement de
l\textquotesingle intelligence en essaim, puis étendue par Heylighen
(2016a, 2016b) aux systèmes numériques contemporains tels que Wikipédia,
le développement open source et les systèmes de contrôle de version.
Ricci et al. (2007) ont proposé le cadre de la \emph{stigmergie
cognitive}, démontrant que des agents rationnels dotés de capacités
avancées peuvent se coordonner via des artefacts environnementaux
manipulables et instrumentables.

L\textquotesingle application de la stigmergie à
l\textquotesingle orchestration d\textquotesingle agents LLM pour la
transformation de code présente un intérêt théorique et pratique
considérable. D\textquotesingle un point de vue théorique, elle propose
de remplacer la communication explicite inter-agents, source de 36,94 \%
des échecs identifiés par MAST (Cemri et al., 2025), par une
coordination médiée par l\textquotesingle environnement.
D\textquotesingle un point de vue pratique, les systèmes de contrôle de
version comme Git fonctionnent déjà comme des médiums stigmergiques où
les commits stimulent les revues, tests et intégrations subséquentes
(Bolici et al., 2016). D\textquotesingle un point de vue managérial, la
traçabilité intrinsèque des artefacts partagés répond aux exigences
croissantes de gouvernance, d\textquotesingle auditabilité et de
conformité réglementaire, notamment celles de l\textquotesingle Article
14 de l\textquotesingle EU AI Act imposant une supervision humaine
effective des systèmes d\textquotesingle IA à haut risque (Fink, 2025).

Toutefois, la littérature académique présente un gap significatif : si
la stigmergie est bien théorisée en biologie et en systèmes distribués
classiques, et si les systèmes multi-agents LLM font
l\textquotesingle objet d\textquotesingle une recherche active, aucun
cadre intégrateur ne propose de combiner coordination stigmergique,
agents LLM et gouvernance organisationnelle dans un framework
généraliste applicable à des workflows complexes, dont la transformation
de code à grande échelle constitue un cas d\textquotesingle étude
privilégié. Les travaux existants traitent ces dimensions de
manière cloisonnée : la littérature sur la stigmergie ignore les LLM, la
littérature sur les agents LLM privilégie les architectures
hiérarchiques, et la littérature en management des systèmes
d\textquotesingle information aborde la gouvernance de
l\textquotesingle IA sans considérer les mécanismes de coordination
émergente. Plus fondamentalement, la théorie de la coordination
(Malone \& Crowston, 1994) identifie la stigmergie comme un mécanisme
organisationnel à part entière, au même titre que les hiérarchies et les
marchés, mais cette perspective n\textquotesingle a jamais été
opérationnalisée pour l\textquotesingle orchestration d\textquotesingle agents IA.

Ce constat motive la problématique suivante :

\begin{quote}
Comment concevoir un framework d\textquotesingle orchestration
multi-agents à coordination stigmergique, généraliste et gouvernable,
capable d\textquotesingle être comparé rigoureusement à des baselines
hiérarchiques reproductibles sur des tâches complexes de coordination
sous contraintes ?
La transformation de code constitue le domaine
d\textquotesingle application principal pour la validation empirique de
ces principes de conception.
\end{quote}

Cette problématique se décline en cinq objectifs de conception,
conformément à la méthodologie Design Science Research adoptée
(Hevner et al., 2004 ; Peffers et al., 2007) :

\begin{itemize}
\tightlist
\item
  OC1, Architecture généraliste : Concevoir un framework
  d\textquotesingle orchestration stigmergique (phéromones numériques,
  règles locales, environnement partagé) applicable à des tâches
  variées de coordination multi-agents, indépendamment du domaine.
\item
  OC2, Coordination émergente : Démontrer que les agents développent
  une spécialisation émergente et une allocation dynamique des tâches
  sans assignation explicite de rôles, conformément aux principes
  des systèmes adaptatifs complexes.
\item
  OC3, Comparaison contrôlée sur la planification sous contraintes :
  Comparer le framework à des baselines centralisées reproductibles sur
  le benchmark TravelPlanner (Xie et al., 2024) avec le même backbone,
  le même protocole d\textquotesingle évaluation et les mêmes
  contraintes dures.
\item
  OC4, Spécialisation à la transformation de code : Adapter le
  framework aux tâches de migration et de modernisation logicielle, et
  démontrer sa compétitivité sur les benchmarks PolyMigration (Amazon,
  2025) et SWE-bench (Jimenez et al., 2024) par rapport aux baselines
  existantes.
\item
  OC5, Gouvernance et auditabilité : Intégrer des mécanismes de
  traçabilité et de conformité réglementaire directement dans
  l\textquotesingle environnement stigmergique, et en valider la
  pertinence par évaluation d\textquotesingle experts.
\end{itemize}

La revue suit une progression continue : d\textquotesingle abord les
fondements de la stigmergie et leur ancrage dans les théories
managériales (coordination organisationnelle, capacités dynamiques,
routines organisationnelles), puis les mécanismes de coordination
computationnelle, les limites des architectures multi-agents, les usages
concrets en migration de code, et enfin les enjeux de gouvernance, de
processus et d\textquotesingle évaluation. L\textquotesingle objectif est
de construire un cadre cohérent, à la fois techniquement crédible et
ancré dans les sciences de gestion, conformément aux exigences du
Design Science Research en systèmes d\textquotesingle information
(Hevner et al., 2004).
